% Ad Familiares 3.12 --- headnote
% ad-familiares-03-12, 4 August 50 BC, Side (south coast of Asia Minor)

\textit{Cicero to Appius Claudius Pulcher, written from
Side on the south coast of Asia Minor on the 4th of
August 50 BC (Perseus dateline: \textit{Scr. Sidae a.
d. iii aut prid. Non. Sext. a. 704 (50)}). Cicero is
sailing west, his year of proconsular command in
Cilicia now over, and lands at Side to find a packet
of letters from his people in Rome. Two pieces of news
shape this letter. The first is the acquittal of
Appius on his second prosecution, the \textit{ambitus}
charge, brought again by Dolabella. The second is
that, in Cicero's absence, his wife Terentia and his
daughter Tullia have concluded Tullia's engagement to
Dolabella himself --- the same Dolabella who has just
twice prosecuted Appius. Cicero is now obliged to
write to Appius from inside a piece of forensic
geometry: he must congratulate Appius on his acquittal
of the very prosecution his daughter's prospective
husband had pressed.}

\textit{Section 1 dispatches the congratulation in
high register. The marvel, Cicero says, is not that
Appius was acquitted --- no one ever doubted that ---
but that in these times, with these manners, no
malice ventured to assail him even in the secrecy of
the ballot. Section 2 then turns to the awkward
matter of the engagement. Cicero asks Appius to step
into his place and feel for him: ``If you find it
easy to say what to say, you needn't make any
allowance for my hesitation.'' He insists that the
match was concluded without his knowledge by his
people in Rome, that the timing was not of his
choosing, and that he has no quarrel with the man
chosen --- but the position is delicate, and he is
plainly sweating it out for the reader to see.}

\textit{Section 3 resolves the difficulty with a
characteristic Ciceronian formula: had he himself
been on the spot, he would have approved the match,
but he would not have settled the timing without
Appius's counsel. Section 4 closes with a small but
affecting domestic scene: as Cicero's ship neared
Side, with Q. Servilius beside him, the letters from
Rome were brought aboard, and Cicero --- seeing that
Servilius too looked shaken --- promised at once that
Appius should expect from him only greater attentions.
The deepening of his old goodwill, with the new
\textit{adfinitas} added on, is the note on which the
letter ends.}
